%--- iliad ---
% separateSolutions: true
% title: Condensation in context
% summary: >-
%   How condensation relates to earlier work in statistics, what its correspondence
%   theorem establishes, and what remains between that theorem and useful
%   translation between the latents of human and AI models.
%--- end ---
\documentclass[10pt]{article}
\IfFileExists{iliad.sty}{\usepackage{iliad}}{\usepackage{../iliad}}
\title{Condensation in context}
\author{Satya Benson}
\date{September 21, 2026}
\begin{document}
\maketitle

\section{Prerequisites}
This reading follows the condensation worksheet. It uses conditional entropy,
conditional independence, and the distinction between an individual latent and
a collection of latents. All variables in the examples have finite alphabets;
logarithms are base two.

\begin{learningoutcomes}
\begin{itemize}
\item Relate condensation to sufficient statistics, common information,
computational mechanics, and graphical models.
\item Separate informational correspondence from finding and using a translation.
\item Identify which additional claims would connect correspondence to AI safety.
\end{itemize}
\end{learningoutcomes}

\section{What is being characterized?}
Fix jointly distributed observables $X_1,\ldots,X_n$. We organize information
into latents $Y_A$, indexed by nonempty subsets $A$ of the observable indices.
The address specifies which observations the latent is assigned to. Each $X_i$
is determined by its contributing collection $(Y_A:i\in A)$.

LVM alone leaves considerable freedom. The additional reconstruction
and Markov conditions constrain where shared information can reside. The
correspondence theorem compares two models satisfying these conditions: their
collections above the same address determine one another, exactly or with a
controlled conditional-entropy error.

This provides a conditional sense of objectivity. Once the observables and the
conditions are fixed, agreement in information content need not come from a
shared vocabulary. It follows from each model's relationship to the observations.
The theorem does not choose the observables. Changing the questions changes the
problem, and may change the information that must correspond.

Nor is each latent automatically a human concept. A model includes residual
information needed to determine observations. The interpretation of its parts
as concepts is an additional connection between the mathematics and the uses of
a latent model.

\section{Connections to prior work}
\subsection{Sufficient statistics}
For a fixed joint law of $(X,V)$, a statistic $T=f(X)$ is sufficient for predicting
$V$ when
\[
 I(X;V\mid T)=0.
\]
Knowing $X$ then gives no predictive information about $V$ beyond $T$.
For finite alphabets, grouping values $x$ with the same conditional distribution
$P(V\mid X=x)$ gives a minimal such statistic, on the support of $X$.
This differs from parametric sufficiency, where the conditional law of
the sample given a statistic must be independent of the unknown parameter.
Both formulations ask which distinctions must be retained for an inference task.

Predictive sufficiency does not say that $T$ determines the realized value of
$V$. Condensation's LVM condition does demand determination, using the whole
contributing collection, including residual variables. A sufficient
statistic can itself be a vector or another structured object.

A sufficient statistic is by construction a function of $X$. A latent need not
be, because reconstruction is asked for only up to a defect: that leaves room
for latents which are parameters of a law rather than functions of the
observations. The additional
question here is how to organize a family of latent parts across many specified
observables. See \href{https://www.stats.org.uk/statistical-inference/Fisher1922.pdf}{Fisher (1922)}
and \href{https://www.stat.cmu.edu/~cshalizi/dm/19/lectures/17/lecture-17.html}{Shalizi's account of predictive sufficiency}.

\subsection{Common information}
Two older questions distinguish requirements that appear together in the
perfect case. The maximal deterministic common part of $X_1,X_2$ is a variable
$K$ recoverable from either observation, through which every other common
function factors. Its entropy is the G\'acs--K\"orner common information.
Wyner's common information instead minimizes $I(X_1,X_2;W)$ over extensions
satisfying $X_1\perp X_2\mid W$. Such a $W$ need not be separately recoverable
from the observations.

If $W$ is both separately recoverable and screens off dependence, then
\[
 I(X_1;X_2)=H(W)+I(X_1;X_2\mid W)=H(W).
\]
Moreover, every common function $K$ is determined by $W$: the two-witness lemma,
with witnesses $X_1,X_2$ and common part $W$, gives $H(K\mid W)=0$.
Thus this special case is closely connected to common information in the sense
of G\'acs and K\"orner.

It is restrictive. If $X$ is a fair bit and $V=X\mathbin{\oplus}N$ for independent
$N\sim\operatorname{Bernoulli}(p)$, $0<p<1/2$, all four pairs have positive
probability. Any common function must therefore be constant, although
$I(X;V)=1-h_2(p)>0$, where
$h_2(p)=-p\log_2 p-(1-p)\log_2(1-p)$. Recoverable common information and
mutual information differ.

Common/private source coding already studies structured organization
of information. Condensation's particular focus is a family of subset-addressed
latents and correspondence under conditions on that family.
See \href{https://cs-web.bu.edu/faculty/gacs/papers/commoninf.pdf}{G\'acs and K\"orner (1973)}
and \href{https://arxiv.org/pdf/1010.3613}{Liu, Xu and Chen's extension of Wyner's common information}.

\subsection{Computational mechanics}
Computational mechanics groups past histories according to their conditional
distributions over futures. Its causal states have predictive optimality,
minimality, and uniqueness properties. This is an important precedent for
characterizing such a construction through informational requirements.

The present setup starts from a family of jointly distributed questions rather
than a distinguished past/future split. It characterizes latent collections
associated with different observable subsets.
See \href{https://arxiv.org/abs/cond-mat/9907176}{Shalizi and Crutchfield (2001)}.

\subsection{Graphical models and the information bottleneck}
The ordered Markov condition in condensation is a graphical-model condition.
A subset-ordered directed graph places larger addresses above smaller ones.
Under the corresponding factorization, the latent joint law is
\[
 P(y)=\prod_{A\ne\varnothing}P(y_A\mid y_{\supsetneq A}).
\]
A factorization alone does not ensure that observations recover the latents,
or that different latent models correspond. Those are additional requirements.
Graphical language also does not make the result a causal-identification theorem.
Eisenstat discusses this connection explicitly in Definition 5.8 and Example 6.3
of \href{https://openreview.net/pdf?id=HwKFJ3odui}{the condensation paper}.

The information bottleneck compresses an input $X$ into $T$ while retaining
information about a specified relevance variable $V$. A standard objective is
\[
 I(X;T)-\beta I(T;V), \qquad T\perp V\mid X.
\]
It provides an established distribution-based approach to useful summaries of
the data.
Our condensation problem instead involves many addressed latents, LVM, and
conditions across observable sets. An optimum of a generic
bottleneck objective does not automatically satisfy these conditions.
See \href{https://arxiv.org/abs/physics/0004057}{Tishby, Pereira and Bialek (1999)}.

\section{What information theory leaves out}
\subsection{An encoding and access to its inverse}
Let $X$ be uniform on $\{1,\ldots,N\}$ and let $Y=\pi(X)$ for a fixed
permutation $\pi$. Then
\[
 H(X\mid Y)=H(Y\mid X)=0.
\]
This is exact correspondence whether $\pi$ is the identity or a large
unstructured lookup table. The entropies do not charge for describing the table,
learning it, storing it, or using it. For a finite table an inverse exists and is
computable; that does not make it available to a particular learner.

For a concrete learning distinction, draw $\pi$ uniformly from all permutations
and reveal $m<N$ distinct input--output pairs. Conditional on those pairs, the
image of a previously unseen input is uniform over the $N-m$ unused outputs. The best success
probability without further information is $1/(N-m)$. For every fixed permutation
the correspondence entropy is zero, yet a learner with few examples cannot
predict an unseen pairing. These are different probability models: one fixes
the permutation; the other describes uncertainty about it.

A separate question concerns computation even when a transformation has a short
known description. Entropic bounds alone give no running-time bound for an
inverse. Establishing computational difficulty requires an additional argument
or assumption; it does not follow from this permutation example.

\subsection{Rare events and average error}
Let $E\sim\operatorname{Bernoulli}(p)$. Define $R=0$ when $E=0$ and, when
$E=1$, let $R$ be uniform over $M$ other values. Since $R$ reveals $E$,
\[
 H(R)=h_2(p)+p\log_2 M.
\]
For $p=10^{-6}$ and $M=2^{10^6}$, the rare branch contributes one full bit.
An independent sample of size $m$ misses that branch with probability
$(1-p)^m$. Large informational contributions can be absent from a small
dataset.

Conversely, suppose a variable $T$ records $E$, and records the value of a fair
bit $B$, independent of $E$, only when $E=0$. Then
\[
 H(B\mid T)=p.
\]
The average uncertainty is small, but conditional on the rare event $E=1$, the
bit is entirely unknown. If that event is a critical decision, its importance
is not measured by its probability alone. Entropy is an average informational
quantity, not a measure of the consequences of errors.

These examples distinguish informational equivalence, statistical access, and
performance on important cases. A useful translation may need guarantees about
all three.

\section{The connection to AI safety}
\begin{enumerate}
\item \textbf{Corresponding information.} Under the reconstruction and Markov
conditions, the theorem gives correspondence between specified collections in
two models of the same observables.
\item \textbf{Concepts of actual agents.} To apply that result, we need to
identify suitable variables in actual agents and establish the relevant
conditions, including how their observables relate. The theorem does not show
that learning procedures produce such models.
\item \textbf{A usable translation.} Informational correspondence could support
identifying which of an AI's latents carries something humans care about. We
must still find
the correspondence with available data and computation, and connect it to the
particular human distinction of interest.
\item \textbf{Interpretation and specification.} A usable translation could help
interpret predictions or express objectives in an AI's world model. This requires
preserving relevant distinctions when circumstances or models change.
\item \textbf{Safer systems.} Better interpretation and specification could improve
oversight and alignment. Sharing concepts does not entail sharing values,
reporting faithfully, or following an intended objective.
\end{enumerate}
The first step is the mathematical result of this lesson. The subsequent steps
state a research programme and the additional work required to connect it to
safety.

\section{Discussion}
\begin{enumerate}
\item \textbf{Choosing the questions.} Consider two agents describing the same
physical system, one using local measurements and another using a different
collection of measurements. What would have to be established before applying
a theorem that assumes the same observables? Distinguish changing value labels
from changing which distinctions an observable makes. Identify one restriction
on observable choice that seems justified for a concrete application.
\item \textbf{What would count as a useful translation?} Suppose two latent tuples $U,V$
satisfy $H(U\mid V)=H(V\mid U)=0$. Specify an additional test
that would make this correspondence useful for interpreting a learned model.
State what access the test requires: paired samples, interventions, model
internals, or a known decoder. Which part of the permutation example would
that access resolve?
\item \textbf{From correspondence to oversight.} Choose one human distinction
that matters for evaluating an AI's behavior. Trace the steps from a
correspondence theorem to an oversight procedure that uses that distinction.
Name a failure that could occur despite exact informational correspondence,
and state what additional evidence would address it.
\end{enumerate}

\section{Further reading}
The linked sources above support different parts
of the discussion. For the mathematical development, start with
\href{https://openreview.net/pdf?id=HwKFJ3odui}{Eisenstat's \emph{Condensation: A Theory of Concepts}},
particularly the exact and quantitative correspondence results.
For the score connection, see
\href{https://satyabenson.com/writing/condensation-conditions-conditioned-score}{\emph{Relating Almost Perfect Condensation Conditions and the Conditioned Score}}.
The score decomposition connects excess conditioned code length to dependence
and recoverability defects; the converse requires attention to which observation
sets the reconstruction assumptions cover.
\end{document}
